% !TeX root = ../Book1.tex
\begin{figure}[htbp]
	\centering
	\begin{tikzpicture}[scale=0.88,font=\small]
		\begin{scope}
			\fill[GeometryBlueLight] (0,0) circle (1.15);
			\draw[GeometryBlue] (0,0) circle (1.15);
			\foreach \a in {0,55,135,210,285}{\draw[GeometryBlue] (0,0)--(\a:1.15);}
			\draw[GeometryBlue] (0.38,0) arc (0:55:0.38);
			\node at (27:0.7) {\(\beta_1\)};
			\fill (0,0) circle (1.3pt);
			\node[align=center] at (0,-1.65) {内点\\\(\sum\beta_j=2\pi\)};
		\end{scope}
		\begin{scope}[xshift=4cm]
			\fill[GeometryBlueLight] (-1.15,0)--(1.15,0) arc (0:180:1.15)--cycle;
			\draw[GeometryBlue] (-1.15,0)--(1.15,0) arc (0:180:1.15);
			\foreach \a in {45,110}{\draw[GeometryBlue] (0,0)--(\a:1.15);}
			\fill (0,0) circle (1.3pt);
			\node at (0,-0.45) {\(\partial M\)};
			\node[align=center] at (0,-1.65) {光滑边界点\\\(\sum\beta_j=\pi\)};
		\end{scope}
		\begin{scope}[xshift=8cm]
			\fill[GeometryBlueLight] (0,0)--(1.15,0) arc (0:125:1.15)--cycle;
			\draw[GeometryBlue] (0,0)--(1.15,0) arc (0:125:1.15)--cycle;
			\foreach \a in {40,83}{\draw[GeometryBlue] (0,0)--(\a:1.15);}
			\draw[GeometryBlue] (0.42,0) arc (0:125:0.42);
			\node at (63:0.72) {\(\beta_v\)};
			\fill (0,0) circle (1.3pt);
			\node[align=center] at (0.1,-1.65) {原有角点\\\(\sum\beta_j=\beta_v\)};
		\end{scope}
	\end{tikzpicture}
	\caption[粘合时的顶点角计数]{各面在同一顶点的切向扇区。内部顶点铺满一周，光滑边界顶点铺满半周，原有角点只铺满原内角。图中的圆弧仅用于标示扇区，不是曲面的分解边。}
	\label{g1:fig:gb-vertex-angles}
\end{figure}
